\begin{python0}
from solutions import *; clear()
\end{python0}
\sectionthree{operator[]}

If \textit{c} is a \textit{C} object, C++ translates \textit{c[n]} like
this:

\textit{c[n]} is the same as \textit{c.operator[](n)}

\textit{operator[]} must be a member function.

\begin{consolethree}[escapeinside=||]
class IntDynArr
{
public:
    IntDynArr(int s)
        : size_(s), p_(new int[s])
    {}
    ~IntDynArr() { delete [] p_; }
    |{\bf int operator[](int n)}| { return p_[n]; }
private:
    int size_;
    int * p_;
};
\end{consolethree}

\begin{consolethree}[escapeinside=||]
int main()
{
    IntDynArr c(100);
    std::cout << c[5] << std::endl;
    return 0;
}
\end{consolethree}

But what about this?

\begin{consolethree}[escapeinside=||]
int main()
{
    IntDynArr c(100);
    c[5] = 23; // ERROR!
    std::cout << c[5] << std::endl;
    return 0;
}
\end{consolethree}

\begin{consolethree}[escapeinside=||]
class IntDynArr
{
public:
    IntDynArr(int s)
        : size_(s), p_(new int[s])
    {}
    ~IntDynArr() { delete [] p_; }
    ...
    |{\bf int \& operator[](int n)}|
    {
        return p_[n];
    }
    ...
private:
    int size_;
    int * p_;
};
\end{consolethree}

Of course you can includes checks like \textit{size >= 0} for
constructor, and argument of \textit{[]} is within bounds:

\begin{consolethree}[escapeinside=||]
...
int & operator[](int n)
{
    if (n >= 0 && n < size)
        return p_[n];
}
...
\end{consolethree}

But what if \textit{n < 0} or \textit{n >= size}?

Later, we will handle such errors with \textbf{exceptions}.

Now include

\begin{consolethree}[escapeinside=||]
...
int & operator[](int n) const
{
    if (n >= 0 && n < size)
        return p_[n];
}
...
\end{consolethree}

What is the point of this?


